\sgasection{5.3}{Use of effectiveness criteria: Noether's theorem}
\nde{In fact, the isomorphisms established in 5.2.8 to 5.2.10
would deserve to be called \og Noether's isomorphism theorems\fg.}
Let \(\mathcal C,\mathcal T\), and (M) be as usual.
Let (N) be a family of morphisms satisfying axioms (a)
and \((\mathrm f_{\mathrm M})\) of 4.7.
Putting together 5.2.7 and 4.7.2, one obtains:
\begin{proposition}{5.3.1}\label{IV.5.3.1}
Let \(G\) be a \(\mathcal C\)-group. Let \(H\) be an
invariant sub-\(\mathcal C\)-group of \(G\), acting
in an (M)-effective manner on \(G\).

For every sub-\(\mathcal C\)-group \(K\) of \(G\)
containing \(H\) and such that the morphism \(K\to G\)
belongs to (N), \(H\) acts on \(K\) in an (M)-effective
manner, and the quotient \(K/H=K'\) is a
sub-\(\mathcal C\)-group of \(G/H=G'\), such that
the morphism \(K'\to G'\) belongs to (N).

The map \(K\mapsto K'\) is a bijection between
the set of sub-\(\mathcal C\)-groups \(K\) of \(G\)
containing \(H\) and such that \(K\to G\) belongs
to (N), and the set of sub-\(\mathcal C\)-groups
\(K'\) of \(G'\) such that \(K'\to G'\) belongs
to (N). The inverse map is \(K'\mapsto K\times_{G'}G\).
Under this correspondence, invariant subgroups of
\(G,G'\) correspond.
% typo?: kept as printed: K rather than K' in the inverse-map formula.
\end{proposition}
\begin{corollary}{5.3.2}\label{IV.5.3.2}
If \(H\to G\) is an element of (N), then
\(\mathcal C\) possesses a final object \(e\),
and the unit section \(e\to G/H\) is an element of (N).
\nde{This follows from the proposition applied to \(K=H\).}
\end{corollary}

\sgasection{6}{Topologies in the category of schemes}
\sgasection{6.1}{The Zariski topology}
This is the topology generated by the following pretopology:
a family of morphisms \(\{S_i\to S\}\) is covering
if each morphism is an open immersion and the union
of the images of the \(S_i\) is all of \(S\).
It is denoted (Zar).
\begin{definition}{6.1.1}\label{IV.6.1.1}
A sheaf for the Zariski topology is also called
a \emph{functor of local nature}: it is a
contravariant functor from \(\Sch\) to \(\Ens\)
such that for every scheme \(S\) and every covering
of \(S\) by open sets \(S_i\), one has an exact diagram:
\[
F(S)\longrightarrow\prod_i F(S_i)
\rightrightarrows\prod_{i,j}F(S_i\cap S_j).
\]
\end{definition}
In particular, a functor of local nature transforms
direct sums into products. Since every representable
functor is a sheaf, this topology is less fine than
the canonical topology.
As regards terminology, whenever we say \og local\fg,
\og locally\fg, without further specification, this
will refer to the Zariski topology, hence have the
customary meaning.

\sgasection{6.2}{A procedure for constructing topologies}
\begin{proposition}{6.2.1}\label{IV.6.2.1}
Let \(\mathcal C\) be a category, \(\mathcal C'\)
a full subcategory, \(P\) a set of families of
morphisms of \(\mathcal C\) with the same target,
stable under base change and composition
(i.e. satisfying axioms (P 1) and (P 2) of 4.2.5),
and \(P'\) a set of families of morphisms of
\(\mathcal C'\) containing the families consisting
of a single identity isomorphism.
Endow \(\mathcal C\) with the topology generated
by \(P,P'\) (cf. 4.2.5.0), and suppose that the
following three conditions hold:

(a) If \(\{S_i\to S\}\in P'\)
(thus \(S_i,S\in\Ob\mathcal C'\)), and if
\(T\to S\) is a morphism of \(\mathcal C'\),
then the fibered products \(S_i\times_S T\)
(in \(\mathcal C\)) exist, and the family
\(\{S_i\times_S T\to T\}\) belongs to \(P'\)
(thus \(S_i\times_S T\in\Ob\mathcal C'\)).
(Remark: this condition implies that \(P'\)
is stable under base change in \(\mathcal C'\),
but is not equivalent to it, for it assumes
moreover that the inclusion functor from
\(\mathcal C'\) to \(\mathcal C\) commutes
with certain fibered products.)

(b) For every \(S\in\Ob\mathcal C\), there
exists \(\{S_i\to S\}\in P\), with
\(S_i\in\Ob\mathcal C'\) for each \(i\).

(c) In the following situation:
\[
\xymatrix{
 & S_{ijk}\ar[d]^{(P')}\\
S_i\ar[d]_{(P')} & S_{ij}\ar[l]_{(P)}\\
S & },
\]
where \(S,S_i,S_{ij},S_{ijk}\in\Ob\mathcal C'\);
\(\{S_i\to S\}\in P'\);
\(\{S_{ij}\to S_i\}\in P\) for each \(i\);
\(\{S_{ijk}\to S_{ij}\}\in P'\) for each \(ij\),
there exists a family \(\{T_n\to S\}\in P'\)
and, for each \(n\), a multi-index \((ijk)\)
and a commutative diagram
\[
\xymatrix{S_{ijk}\ar[dr] & T_n\ar[l]\ar[dl]\\ & S }.
\]
Then, for a sieve \(R\) of \(S\in\Ob\mathcal C\)
to be covering, it is necessary and sufficient
that there exist a composite family
\[
\xymatrix{
R\ar[d] & S_{ij}\ar@{-->}[l]\ar[d]^{(P')}\\
S & S_i\ar[l]_{(P)} },
\]
where \(S_i,S_{ij}\in\Ob\mathcal C'\),
\(\{S_i\to S\}\in P\), \(\{S_{ij}\to S_i\}\in P'\)
for each \(i\), and that the resulting morphisms
\(S_{ij}\to S\) factor through \(R\)
(in other words, that the sieve generated by
this composite family be contained in \(R\)).
\end{proposition}
\begin{proof}
Since families that are elements of \(P,P'\)
are covering, a composite family of such
families will also be covering (C 2);
hence a sieve of the indicated form will
be covering, since it contains a covering sieve.

Conversely, it suffices to see that the sieves
of this form do form a topology, i.e. it
suffices to verify axioms (T 1) to (T 4) of 4.2.1.

\emph{Axiom (T 4).} Let \(S\in\Ob\mathcal C\).
By (b), there exists a family \(\{S_i\to S\}\in P\),
with \(S_i\in\Ob\mathcal C'\).
The families \(\{S_i\xrightarrow{\mathrm{id}_{S_i}}S_i\}\)
are elements of \(P'\) by hypothesis.
The sieve \(S\) of \(S\) is therefore of the desired form:
\[
\xymatrix{S_i\ar[r]\ar[d]_{(P')\;\mathrm{id}} & S\ar[d]^{\mathrm{id}_S}\\
S_i\ar[r]_{(P)} & S }.
\]

\emph{Axiom (T 3).} Evident.

\emph{Axiom (T 2).} Let \(R\) be a sieve of
\(S\) of the desired form, and let \(C\) be
a sieve\nde{The original has been corrected here.}
of \(S\) such that, for every \(T\in\Ob\mathcal C\)
and every morphism \(T\to S\) factoring through
\(R\), the sieve \(C\times_S T\) of \(T\)
is of the desired form.
Then, since \(S_{ij}\to S\) factors through \(R\),
the sieve \(C_{ij}=C\times_S S_{ij}\) of \(S_{ij}\):
\[
\xymatrix{
 & S_{ij}\ar[d]^{(P')}\ar[ddl] & C_{ij}\ar@{^{(}->}[l]\ar[dd]\\
 & S_i\ar[d]^{(P)} & \\
R\ar@{^{(}->}[r] & S & C\ar@{^{(}->}[l] }
\]
is of the desired form; thus, for every \(ij\),
one has a diagram of the form:
\[
\xymatrix{
S_{ijkl}\ar[d]_{(P')}\ar[dddr]\\
S_{ijk}\ar[d]_{(P)}\\
S_{ij} & \\
 & C_{ij}\ar@{^{(}->}[ul] }.
\]
One has therefore proved that there exists
a composite family
\[
S_{ijkl}\xrightarrow{(P')}S_{ijk}
\xrightarrow{(P)}S_{ij}\xrightarrow{(P')}S_i
\xrightarrow{(P)}S
\]
belonging to \(P\circ P'\circ P\circ P'\),
factoring through \(C\), and in which all
objects other than \(S\) are in \(\mathcal C'\).
Applying condition (c) to each family
\(\{S_{ijkl}\to S_i\}\), one deduces that
for each \(i\) there exists a family
\(\{T_{in}\to S_i\}\in P'\) such that
\(T_{in}\to S\) factors through one of the
\(S_{ijkl}\), hence through \(C\):
\[
\xymatrix{
T_{in}\ar[r]\ar[d]_{(P')} & S_{ijkl}\ar[dd]\\
S_i\ar[d]_{(P)} & \\
S & C\ar@{^{(}->}[l] }.
\]
The sieve \(C\) of \(S\) is therefore of the
desired form, which completes the verification.

\emph{Axiom (T 1).} Let \(R\) be a sieve
of \(S\) of the given form, and \(T\to S\)
a morphism of \(\mathcal C\).
Let us show that the sieve \(R\times_S T\)
of \(T\) is of the desired form.
\[
\xymatrix@C=2.1em{
 & S_{ij}\ar[d]^{(P')}\ar[ddl] & & U_{ikj}\ar[ll]\ar[d]^{(P')}\\
 & S_i\ar[d]^{(P)} & T_i\ar[l]\ar[d]^{(P)}
 & U_{ik}\ar[l]_{(P)}\ar[dl]^{(P)}\\
R\ar@{^{(}->}[r] & S & T\ar[l] & }.
\]
Let \(T_i=S_i\times_S T\).
The family \(\{T_i\to T\}\) belongs to \(P\)
(by (P 1)). Applying (b), one constructs
\(\{U_{ik}\to T_i\}\in P\), with
\(U_{ik}\in\Ob\mathcal C'\).
By hypothesis (condition (P 2) on \(P\)),
one has \(\{U_{ik}\to T\}\in P\).
By (a), \(U_{ik}\times_{S_i}S_{ij}=U_{ikj}\)
is an object of \(\mathcal C'\), and for
every \(ik\), \(\{U_{ikj}\to U_{ik}\}\in P'\).
Then the commutative diagram below
\[
\xymatrix{
R\ar[d] & & U_{ikj}\ar[ll]\ar[d]^{(P')}\\
S & T\ar[l] & U_{ik}\ar[l]_{(P)} }
\]
shows that the morphisms \(U_{ijk}\to T\)
factor through the sieve \(T\times_S R\)
of \(T\), which is therefore of the desired form,
and this completes the proof.
% typo?: kept as printed: U_ijk here, U_ikj in the diagram.
\end{proof}

\begin{corollary}{6.2.2}\label{IV.6.2.2}
If \(S\in\Ob\mathcal C'\) and \(R\) is a
sieve of \(S\), \(R\) is covering if and
only if there exists a family
\(\{T_i\to S\}\in P'\) factoring through \(R\).
\end{corollary}
\begin{proof}
Indeed, such a sieve is covering.
On the other hand, it suffices to apply
(c), taking the family \(\{S_i\to S\}\)
to consist of the identity isomorphism
of \(S\), to deduce from the proposition
that a covering sieve is of the indicated form.
\end{proof}
\begin{corollary}{6.2.3}\label{IV.6.2.3}
For a presheaf \(F\) on \(\mathcal C\)
to be separated (resp. a sheaf), it is
necessary and sufficient that the morphism
\[
F(S)\longrightarrow\prod_i F(S_i)
\]
be injective (resp. that the diagram
\[
F(S)\longrightarrow\prod_i F(S_i)
\rightrightarrows\prod_{i,j}F(S_i\times_S S_j)
\]
be exact) in the following two cases:

(i) \(\{S_i\to S\}\in P\),

(ii) \(S,S_i\in\Ob\mathcal C'\);
\(\{S_i\to S\}\in P'\).
\end{corollary}
\begin{proof}
Indeed, the conditions are necessary,
for the families in question are covering.
If \(C\) is the sieve of \(S\) which is
the image of a family of morphisms
\(\{S_{ij}\xrightarrow{(P')}S_i\xrightarrow{(P)}S\}\),
an easy diagram chase shows that the
conditions of the corollary imply that
\(\Hom(S,F)\xrightarrow{g}\Hom(C,F)\)
is injective (resp. bijective).
But every refinement \(R\) of \(S\)
contains a sieve \(C\) of the above type,
and one has a commutative diagram
\[
\xymatrix{
\Hom(S,F)\ar[rr]^f\ar[dr]_g & & \Hom(R,F)\ar[dl]^h\\
 & \Hom(C,F) & }.
\]
One knows that \(g\) is injective,
hence \(f\) is too. Thus \(F\) is separated.
But \(R\) is a refinement of \(C\),
hence \(h\) is also injective.
If \(g\) is bijective, then \(f\)
is too, hence \(F\) is a sheaf.
\end{proof}
\begin{remark}{6.2.4}\label{IV.6.2.4}
The preceding corollary does not follow
from 4.3.5, since \(P'\) is not stable
under base extension.
\end{remark}
\begin{remark}{6.2.5}\label{IV.6.2.5}
Condition (c) holds in particular in the case where

(i) \(P'\) is stable under composition.

(ii) If \(\{S_j\to S\}\) is a family
of morphisms of \(\mathcal C'\), an
element of \(P\), there exists a
subfamily of it which is an element of \(P'\).
\end{remark}

\sgasection{6.3}{Application to the category of schemes}
Take for \(\mathcal C\) the category of schemes,
for \(\mathcal C'\) the full subcategory
consisting of affine schemes, and for
\(P\) the set of surjective families
of open immersions.
One will consider several sets \(P'\):

\(P'_1\): finite surjective families
consisting of flat morphisms.

\(P'_2\): finite surjective families
consisting of flat morphisms of finite
presentation which are quasi-finite.
\nde{Let \(P''_2\) be the set of finite
surjective families consisting of flat
morphisms of \(\mathcal C'\) of finite
presentation. By Proposition 6.3.1 below,
the topology \(\mathcal T_2\) generated
by \(P,P'_2\) coincides with that
generated by \(P,P''_2\).
This follows from the results of
EGA IV$_4$, \(\S\) 17.16 on quasi-sections;
see the proof of 6.3.1.
Let us cite here the following particular
case of EGA IV$_4$, 17.16.2:
let \(S\) be affine and \(f:X\to S\)
a surjective flat morphism locally of
finite presentation; then there exists
a faithfully flat morphism \(S'\to S\)
of finite presentation which is quasi-finite,
with \(S'\) affine, and an \(S\)-morphism \(S'\to X\).}

\(P'_3\): finite surjective families
consisting of étale morphisms.

\(P'_4\): finite surjective families
consisting of finite étale morphisms.

For each of these sets \(P'_i\), except
\(P'_4\), the conditions of Proposition
6.2.1 hold ((c) by 6.2.5, for, since
an affine scheme is quasi-compact,
every family of morphisms of \(\mathcal C'\)
which is an element of \(P\) contains
a finite subfamily which is also in
\(P\), hence in \(P'_i\) for \(i=1,2,3\)).
The topology \(\mathcal T_i\) generated
by \(P,P'_i\) is denoted and called as follows:
\footnotemark[57]
\[
\begin{aligned}
\mathcal T_1&=(\mathrm{fpqc})=
&&\text{faithfully flat quasi-compact topology}.\\
\mathcal T_2&=(\mathrm{fppf})=
&&\text{faithfully flat (locally) finite presentation topology}.\\
\mathcal T_3&=(\mathrm{\acute et})=
&&\text{étale topology}.\\
\mathcal T_4&=(\mathrm{\acute etf})=
&&\text{finite étale topology}.
\end{aligned}
\]
Since \(P'_1\supset P'_2\supset P'_3\supset P'_4\), one has
\[
(\mathrm{fpqc})\geq(\mathrm{fppf})
\geq(\mathrm{\acute et})\geq(\mathrm{\acute etf})\geq(\mathrm{Zar}).
\]
\begin{proposition}{6.3.1}\label{IV.6.3.1}
(i) For the sieve \(R\) of \(S\) to
be covering for \(\mathcal T_i\),
\(1\leq i\leq3\), it is necessary and
sufficient that there exist a covering
\((S_p)\) of \(S\) by affine open sets
and, for each \(p\), a family
\(\{S_{pq}\to S_p\}\) which is an
element of \(P'_i\), the \(S_{pq}\)
being affine, such that each morphism
\(\{S_{pq}\to S\}\) factors through \(R\).
\nde{By hypothesis, each family
\(\{S_{pq}\to S\}\in P'_i\) is finite,
hence \(S'_p=\coprod_q S_{pq}\) is
affine, and the family can therefore
be replaced by the morphism \(S'_p\to S_p\),
which still belongs to \(P'_i\).}

(ii) For a presheaf \(F\) on \(\Sch\)
to be a sheaf for (fpqc) (resp.
(fppf), (ét), (étf)), it is necessary
and sufficient that

a) \(F\) be a sheaf for (Zar),
i.e. a functor of local nature.

b) For every faithfully flat (resp.
faithfully flat of finite presentation
and quasi-finite, resp. surjective
étale, resp. surjective finite étale)
morphism \(T\to S\), where \(T,S\)
are affine, one have an exact diagram:
\[
F(S)\longrightarrow F(T)\rightrightarrows F(T\times_S T).
\]

(iii) The topologies \(\mathcal T_i\),
\(i=1,2,3,4\), are less fine than
the canonical topology.

(iv) Every surjective family consisting
of flat and open (resp. flat and
locally of finite presentation,
resp. étale, resp. finite étale)
morphisms is covering for (fpqc)
(resp. (fppf), resp. (ét), resp. (étf)).
